\documentclass[11pt]{article}
\usepackage[margin=1in]{geometry}
\usepackage{amsmath,amssymb,amsthm}
\usepackage{enumitem}
\usepackage[hidelinks]{hyperref}

\newtheorem{definition}{Definition}
\newtheorem{lemma}{Lemma}
\newtheorem{theorem}{Theorem}
\newcommand{\R}{\mathbb{R}}
\newcommand{\N}{\mathbb{N}}

\title{A Relaxation Algorithm for Quadratic Energy\\
\large A source-only Lea dependency and retry experiment}
\author{}
\date{}

\begin{document}
\maketitle

We study a scalar relaxation scheme for minimizing a shifted quadratic.
The algebraic identities are exact; convergence requires a genuine
restriction on the step size. The proofs deliberately use earlier results,
so their dependencies are mathematical, not merely a list of related claims.

There are nine Lea targets below: two definitions, six dependent results,
and one independent control theorem. Ordinary LaTeX packages suffice.
All baseline statements are true. Use only this document and Lea's
Define, Formalize, and Re-formalize actions; no Lean editing is required.

\section{The model}

\begin{definition}[Quadratic energy]\label{def:descent-energy}
% lea: define label=descent_energy
% lea: context={Use explicit real arguments a and x. Define the computational real-valued function descent_energy a x = (1/2)*(x-a)^2; do not introduce it as an axiom.}
For $a,x\in\R$, define the energy
\[
  E_a(x)=\frac12(x-a)^2.
\]
\end{definition}

\begin{definition}[Relaxation step]\label{def:descent-step}
% lea: define label=descent_step
% lea: context={Use explicit real arguments a, eta, x, in that order. Define descent_step a eta x = x - eta*(x-a). No step-size assumption is part of the definition.}
For $a,\eta,x\in\R$, define
\[
  T_{a,\eta}(x)=x-\eta(x-a).
\]
\end{definition}

\begin{lemma}[Centered update]\label{lem:descent-error}
% lea: formalize label=descent_error uses={descent_step}
% lea: context={Use the existing descent_step definition. Follow the adjacent algebraic proof and retain the unrestricted real parameter eta.}
For all $a,\eta,x\in\R$,
\[
  T_{a,\eta}(x)-a=(1-\eta)(x-a).
\]
\end{lemma}
\begin{proof}
Expand $T_{a,\eta}(x)-a=x-\eta(x-a)-a$.
Collect the two occurrences of $x-a$ to obtain
$(x-a)-\eta(x-a)=(1-\eta)(x-a)$.
\end{proof}

\section{Energy dissipation}

\begin{theorem}[Exact one-step energy identity]\label{thm:descent-energy-step}
% lea: formalize label=descent_energy_step uses={descent_energy,descent_error}
% lea: context={Apply the checked descent_error lemma before simplifying the energy. Reuse the existing energy and step definitions. Preserve the associated proof's dependence on the centered-update lemma rather than independently expanding the whole target.}
For all $a,\eta,x\in\R$,
\[
  E_a\bigl(T_{a,\eta}(x)\bigr)=(1-\eta)^2 E_a(x).
\]
\end{theorem}
\begin{proof}
By the centered-update lemma,
\[
  T_{a,\eta}(x)-a=(1-\eta)(x-a).
\]
Substitute this equality into the definition of $E_a$ and square:
\[
  \frac12\bigl((1-\eta)(x-a)\bigr)^2
  =(1-\eta)^2\frac12(x-a)^2.
\]
\end{proof}

\begin{theorem}[Strict energy decrease]\label{thm:descent-strict}
% lea: formalize label=descent_strict uses={descent_energy,descent_energy_step}
% lea: context={Apply the existing descent_energy_step theorem. Keep both step-size inequalities and x not equal to a. Prove positivity of the original energy and that the multiplicative factor is strictly below one.}
For all $a,\eta,x\in\R$, if $0<\eta<2$ and $x\ne a$, then
\[
  E_a\bigl(T_{a,\eta}(x)\bigr)<E_a(x).
\]
\end{theorem}
\begin{proof}
The assumption $0<\eta<2$ gives $-1<1-\eta<1$, hence
$(1-\eta)^2<1$. Since $x\ne a$, the energy $E_a(x)$ is
strictly positive. Multiply $(1-\eta)^2<1$ by this positive
energy and apply the exact one-step energy identity.
\end{proof}

\begin{theorem}[Exact geometric energy law]\label{thm:descent-geometric}
% lea: formalize label=descent_geometric uses={descent_energy,descent_step,descent_energy_step}
% lea: context={Use u : Nat -> Real and the recurrence as a hypothesis for every n. Induct on n and apply the checked descent_energy_step theorem at the successor step. The identity is valid for every real eta, including unstable step sizes.}
Let $a,\eta\in\R$ and let $u:\N\to\R$ satisfy
\[
  u(n+1)=T_{a,\eta}\bigl(u(n)\bigr)
  \qquad\text{for every }n\in\N.
\]
Then for every $n\in\N$,
\[
  E_a\bigl(u(n)\bigr)=\bigl((1-\eta)^2\bigr)^n E_a\bigl(u(0)\bigr).
\]
\end{theorem}
\begin{proof}
Induct on $n$. The case $n=0$ follows because a zeroth power is one.
For the successor step, use the recurrence and the one-step energy identity:
\[
  E_a\bigl(u(n+1)\bigr)=(1-\eta)^2 E_a\bigl(u(n)\bigr).
\]
Insert the induction hypothesis and use the successor law for powers.
\end{proof}

\begin{theorem}[Convergence of the iterates]\label{thm:descent-converges}
% lea: formalize label=descent_converges uses={descent_step,descent_error}
% lea: context={Use u : Nat -> Real and the stated recurrence as a hypothesis. Follow the proof: first derive the centered-iterate formula by induction using the checked descent_error lemma, then use the standard real geometric-sequence limit for absolute ratio less than one. Express convergence with Tendsto atTop (nhds a).}
Let $a,\eta\in\R$ satisfy $0<\eta<2$, and let $u:\N\to\R$
satisfy
\[
  u(n+1)=T_{a,\eta}\bigl(u(n)\bigr)
  \qquad\text{for every }n\in\N.
\]
Then $u(n)\to a$ as $n\to\infty$.
\end{theorem}
\begin{proof}
Induction using the centered-update lemma gives the intermediate formula
\[
  u(n)-a=(1-\eta)^n\bigl(u(0)-a\bigr).
\]
The base case is immediate. In the successor case, substitute the recurrence,
apply the centered-update lemma, and then use the induction hypothesis.
Since $|1-\eta|<1$, the real geometric sequence $(1-\eta)^n$ tends
to zero. Multiplication by the fixed real number $u(0)-a$ preserves
that limit. Adding $a$ proves the claim.
\end{proof}

\begin{lemma}[Unit-displacement calibration]\label{lem:descent-calibration}
% lea: formalize label=descent_calibration uses={descent_energy}
% lea: context={Unfold the existing descent_energy definition and simplify at x=a+1. This is a check of the actual definition body, not an additional axiom.}
For every $a\in\R$,
\[
  E_a(a+1)=\frac12.
\]
\end{lemma}
\begin{proof}
Substitute $x=a+1$ into $E_a(x)=\frac12(x-a)^2$.
The difference $(a+1)-a$ is one, so the energy is $\frac12$.
\end{proof}

\section{Independent control}

\begin{theorem}[Parallelogram identity]\label{thm:independent-square}
% lea: formalize label=independent_square
% lea: context={This is independent of the relaxation model. Prove it by polynomial algebra over Real without importing the project's energy or relaxation results.}
For all $x,y\in\R$,
\[
  (x+y)^2+(x-y)^2=2x^2+2y^2.
\]
\end{theorem}
\begin{proof}
Expand both squares. The mixed terms cancel, leaving $2x^2+2y^2$.
\end{proof}

\section{Source-only experiments}

The following are experiment instructions, not extra assumptions or
replacement hypotheses for the baseline theorems.
For each destructive source experiment, start from a fresh copy of the
fully formalized baseline, or restore the source changes before the next
experiment. Keep all Lea labels unchanged.

\subsection*{1. Build the dependency chain}
\begin{enumerate}[leftmargin=*]
\item Define \texttt{descent\_energy} and \texttt{descent\_step}.
\item Formalize \texttt{descent\_error}, then
\texttt{descent\_energy\_step}.
\item Formalize \texttt{descent\_strict},
\texttt{descent\_geometric}, and \texttt{descent\_calibration}.
\item Formalize \texttt{independent\_square}.
\item Finally formalize \texttt{descent\_converges}; it adds the
analysis-library and local induction-helper case.
\end{enumerate}
Inspect the generated proof summaries or read-only code to confirm that the
named dependencies are actually used. The attached proofs and metadata request
reuse, but a model can still choose another proof. A direct independent proof
has not demonstrated the intended dependency edge.

\subsection*{2. Retry an upstream result with completed consumers}
After the chain is complete, invoke \textbf{Re-formalize} on
\texttt{descent\_error}, then on \texttt{descent\_energy\_step},
without changing their source. Use the explicit retry action rather than
ordinary chat.

Check that Lea admits the retry, retains the target's conversation, reuses
existing definitions/helpers, and preserves downstream proof source.
After actual rechecking, the unchanged mathematics should remain usable.
Reload once during a retry and check that progress reattaches correctly.

For a proof-only source revision, change the proof of the centered-update
lemma to the equivalent argument:
\begin{quote}
Write $x=a+(x-a)$ in the definition of $T_{a,\eta}(x)$.
Subtract $a$ and collect the coefficient of $x-a$.
\end{quote}
Then re-formalize that lemma again. Its mathematical interface is unchanged.

\subsection*{3. Change a definition body and expose a real stale consumer}
In the energy definition only, replace
\[
  E_a(x)=\frac12(x-a)^2
  \quad\text{with}\quad
  E_a(x)=(x-a)^2.
\]
Also change that definition's context comment so its stated body matches
the new formula. Leave the other statements and proofs unchanged.
Invoke \textbf{Regenerate definition} or \textbf{Re-formalize}, as offered.

The one-step and geometric energy identities are still mathematically true:
the common scaling cancels. The calibration lemma, however, is now false,
because the new value at $a+1$ is $1$, not $\frac12$.
It must not remain certified current against the changed definition.
An old proof displayed in history is not a new passing check.

Re-formalize the calibration lemma while its source still says $\frac12$.
Expect an honest obstruction or verified disproof, not a silent change
of its conclusion. Then change its statement and proof to value $1$,
keep its Lea label, and re-formalize it to repair the source.
The independent parallelogram theorem should retain its source and
completion lineage; it has no dependency on this definition.

\subsection*{4. Test a realistic missing stability restriction}
On a fresh baseline copy, change only the strict-decrease theorem's
hypothesis from $0<\eta<2$ to $0<\eta<3$.
Leave its proof as written and invoke \textbf{Re-formalize}.

The revised claim is false. At $a=0$, $x=1$, $\eta=5/2$,
\[
  T_{0,5/2}(1)=-\frac32,\qquad
  E_0(1)=\frac12,\qquad E_0(-3/2)=\frac98.
\]
Lea should identify the faulty hypothesis or produce a verified
counterexample. It must not silently restore $\eta<2$ and call the
revised claim proved. Other valid results and their source must survive.

Restore the original hypothesis and invoke the explicit retry action.
Check that the corrected source can complete in the same target history.

\subsection*{5. Observe interruption and queue behavior}
During the convergence proof, use the available Stop or Pause action
after useful progress appears. Reload and inspect retained progress.
Continue explicitly using the offered action; do not edit generated Lean.

If a second explicit retry can be submitted while the first is active,
submit it for another target in this project. Check for queued activity
and serial execution. If the UI prevents a second submission, record this
queue case as untested. If the first proof finishes before interruption,
record the interruption case as inconclusive.

\paragraph{Coverage boundary.}
This document naturally exercises dependency reuse, source-driven retries,
definition-body freshness, consumer invalidation, retained progress, and
honest failure on a false revision. It does not force a particular Lean file
layout. If Lea places supporting declarations and a retry target in the
same module, those retries also exercise shared-file preservation.
If it generates separate modules, do not count the run as proof of
same-file preservation. No manual Lean edits are needed in either case.

\end{document}
